\documentclass[10pt]{beamer}

\usetheme{metropolis}
\definecolor{WiLabRed}{RGB}{197,18,48}
\setbeamercolor{frametitle}{fg=white,bg=WiLabRed}
\setbeamercolor{progress bar}{fg=WiLabRed!90}
\setbeamercolor{title separator}{fg=WiLabRed!90}
\setbeamercolor{progress bar in section page}{fg=WiLabRed!90}
\setbeamercolor{background canvas}{bg=white}
\setbeamercolor{alerted text}{fg=WiLabRed!90}

\usepackage{appendixnumberbeamer}

\usepackage{booktabs}
\usepackage[scale=2]{ccicons}

\usepackage{pgfplots}
\usepgfplotslibrary{dateplot}

\usepackage{xspace}
\newcommand{\themename}{\textbf{\textsc{metropolis}}\xspace}

\usepackage{marvosym}
%\usepackage{subfig}
\usepackage{graphicx}\graphicspath{{images/}}
\usepackage{subcaption}
\usepackage[framed]{mcode}
\usepackage{listings}

\title{Lecture 14: Error Vector Magnitude Measurement, Phase Ambiguity}
\subtitle{\textit{Software Defined Radio for Engineers} (Collins~\textit{et al.}), \textsection{7.2.4} \& \textsection{7.3}}
\date{}
\author{\textbf{Alexander M. Wyglinski, Ph.D.}}
\institute{ \vspace*{1in}\hfill\includegraphics[height=1.125cm]{wilab_logo-A70916.eps} \qquad \includegraphics[height=1.125cm]{WPI_Inst_Prim_FulClr.eps}}
% \titlegraphic{\hfill\includegraphics[height=1.5cm]{logo.pdf}}

% Foot for all slides
\setbeamertemplate{frame footer}{\tiny \copyright~2018 by Alexander Wyglinski. This work is licensed under the Creative Commons Attribution-ShareAlike 4.0 International License. To view a copy of this license, visit http://creativecommons.org/licenses/by-sa/4.0/.}

\begin{document}

%\captionsetup[subfigure]{labelformat=empty}

%%%%%%%%%%%%%%%%%%%%%%%%%%%%%%%%%%%%%%%%%%%%%%%%%%%%%%%%%%

\maketitle




%%%%%%%%%%%%%%%%%%%%%%%%%%%%%%%%%%%%%%%%%%%%%%%%%%%%%%%%%%

\frame
{
  \frametitle{Error Vector Magnitude Measurement}

    \begin{itemize}
        \item EVM provides useful measure of phase noise in recovered signal
        \begin{itemize}
         \item Helps understand algorithmic performance of Systems
         \item Measures residual error of constellation with respect to reference position
        \end{itemize}
        \item Following mathematical expression used to calculate EVM in percent RMS:
        \begin{equation}
         \mathrm{EVM}_{\mathrm{RMS}}=100\cdot\left(\sqrt{\sum\limits_{k=0}^{N-1}e_{\mathrm{const}}(k)}\sum\limits_{k=0}^{N-1}((\Re(\bar{y}(k)))^2+(\Im(\bar{y}(k)))^2)\right)
        \end{equation}
	and:
	\begin{equation}
	 e_{\mathrm{const}}(k)=(\Re(y(k))-\Re(\bar{y}(k)))^2+(\Im(y(k))-\Im(\bar{y}(k)))^2
	\end{equation}
	where $\bar{y}(k)$ is the reference symbol for $y(k)$
    \end{itemize}


}

%%%%%%%%%%%%%%%%%%%%%%%%%%%%%%%%%%%%%%%%%%%%%%%%%%%%%%%%%%

\frame
{
  \frametitle{Error Vector Magnitude Measurement}
  
  \begin{itemize}
   \item EVM is a measure of dispersiveness of received signal $\rightarrow$ Lower EVM is better
   \item Sometimes EVM is represented in decibels:
   \begin{equation}
    \mathrm{EVM}_{\mathrm{dB}}=20\log_{10}\left(\frac{\mathrm{EVM}_{\mathrm{RMS}}}{100}\right)
   \end{equation}
  \end{itemize}

  
}

%%%%%%%%%%%%%%%%%%%%%%%%%%%%%%%%%%%%%%%%%%%%%%%%%%%%%%%%%%

\frame
{
  \frametitle{Phase Ambiguity}
  
  \begin{itemize}
   \item Occurs when FFC synchronizer is blind to true orientation of transmitted signal
   \item Symmetric modulation schemes may have a number of convergent orientations for its signal constellation 
   \item There exists several approaches to help achieve the correct signal constellation] orientation
   \begin{itemize}
    \item Codewords
    \item Differential Encoding
    \item Equalizers
   \end{itemize}
  \end{itemize}

  
}

%%%%%%%%%%%%%%%%%%%%%%%%%%%%%%%%%%%%%%%%%%%%%%%%%%%%%%%%%%

\frame
{
  \frametitle{Differential Encoding}
  
  \begin{itemize}
   \item To correct for phase ambiguity, differentially encode the source bits
   \begin{itemize}
    \item Makes true data dependent on difference between successive bits instead of just the received bits themselves
   \end{itemize}
    \item Encoding performed using modulo-2 addition:
    \begin{equation}
     b_t(n)=b_t(n-1)\oplus{b(n)}
    \end{equation}
    while decoding is performed using the following expression:
    \begin{equation}
     b(n)=b_t(n)\oplus{b_t(n-1)}
    \end{equation}
    where $b_t(n)$ are the transmitted bits while $b(n)$ are the untransmitted bits
  \end{itemize}

  
}

%%%%%%%%%%%%%%%%%%%%%%%%%%%%%%%%%%%%%%%%%%%%%%%%%%%%%%%%%%




\end{document}
